\section{Conclusiones}

\begin{frame}{Conclusiones Principales}
  \begin{block}{Hallazgos Clave}
    \begin{enumerate}
      \item \textbf{Mayor Supervivencia:} La tabla construida muestra tasas de mortalidad sistemáticamente menores a MI2014 en todos los rangos etarios analizados.
      \item \textbf{Impacto Financiero Significativo:}
        \begin{itemize}
          \item Aumento de 20\% a 186\% en costos de productos de supervivencia
          \item Diferenciales máximos en rentas diferidas para edades jóvenes
          \item Convergencia en edades avanzadas ($>$ 85 años)
        \end{itemize}
      \item \textbf{Riesgo de Insuficiencia:} El uso de MI2014 podría subestimar reservas técnicas necesarias para esta población específica.
    \end{enumerate}
  \end{block}
\end{frame}

\begin{frame}{Implicancias Actuariales}
  \textbf{Para Aseguradoras:}
  \begin{itemize}
    \item Necesidad de ajustar tarifas de productos de supervivencia
    \item Revisión de reservas técnicas para población con características similares
    \item Consideración de tablas específicas por segmento poblacional
  \end{itemize}
  
  \vspace{0.5cm}
  
  \textbf{Para Reguladores:}
  \begin{itemize}
    \item Importancia de tablas actualizadas y específicas
    \item Evaluación periódica de experiencia vs. tablas regulatorias
    \item Requerimientos de suficiencia patrimonial adecuados
  \end{itemize}
\end{frame}

\begin{frame}{Metodología Exitosa}
  \begin{block}{Aspectos Validados}
    \begin{itemize}
      \item \textbf{Graduación Whittaker-Henderson:} Todos los tests estadísticos validaron el ajuste (Chi-cuadrado, KS, Rachas, Signos)
      \item \textbf{Extrapolación:} Métodos de Oppermann y Heligman-Pollard apropiados para cabeza y cola
      \item \textbf{Consistencia:} Tabla resultante muestra comportamiento coherente en todo el rango etario
    \end{itemize}
  \end{block}
\end{frame}

\begin{frame}{Limitaciones del Estudio}
  \begin{itemize}
    \item \textbf{Exposición Limitada:} Pocos datos en edades extremas (0-19 y 85+), requiriendo extrapolación matemática
    \item \textbf{Antigüedad de Datos:} Período 2000-2012 puede no capturar mejoras recientes en mortalidad (avances médicos, cambios demográficos)
    \item \textbf{Especificidad Poblacional:} Resultados válidos solo para población con características similares (Origen 2, Masculino)
    \item \textbf{Supuesto de Suavidad:} Método Whittaker-Henderson asume suavidad que podría enmascarar discontinuidades reales
    \item \textbf{Tasa de Interés Fija:} Análisis con $i = 5\%$ constante; variaciones afectarían valores actuales
  \end{itemize}
\end{frame}

\begin{frame}{Recomendaciones}
  \begin{enumerate}
    \item \textbf{Actualización Periódica:} Revisar tabla cada 3-5 años con datos más recientes
    \item \textbf{Segmentación:} Considerar construcción de tablas por sub-poblaciones (región, nivel socioeconómico, etc.)
    \item \textbf{Análisis de Sensibilidad:} Evaluar impacto de variaciones en tasa de interés técnica
    \item \textbf{Validación Continua:} Comparar experiencia real vs. tabla periódicamente
    \item \textbf{Mejora de Datos:} Incrementar exposición en edades extremas para reducir dependencia de extrapolación
  \end{enumerate}
\end{frame}

\begin{frame}
  \centering
  \Huge \textcolor{celesteoscuro}{\textbf{¡Gracias!}}
  
  \vspace{1cm}
  
  \Large \textcolor{celesteprincipal}{Preguntas y Comentarios}
  
  \vspace{1.5cm}
  
  \normalsize
  \textit{Juan Pablo Cuevas \& Esteban Román Soto}
  
  \vspace{0.3cm}
  
  \small
  \textcolor{grisoscuro}{Pontificia Universidad Católica de Chile}
  
  \textcolor{grisoscuro}{Diciembre 2025}
\end{frame}
